\section{Introduction}
\label{sec:introduction}

Object navigation asks an embodied agent to find an instance of a requested category from egocentric observations. Progress in semantic mapping and foundation-model reasoning has made this problem increasingly tractable in static environments \citep{chaplot2020objectnav,huang2023vlmaps,gu2024conceptgraphs}. A household robot deployed over weeks, however, rarely acts in the scene that it originally mapped. A mug last observed beside the sink may have moved to a desk after breakfast; keys may alternate between a shelf and a bag; and a rarely moved appliance may remain exactly where it was. In each case, the agent has a valid historical observation but an uncertain current navigation goal.

This distinction exposes a limitation of conventional embodied memory. Most spatial memories answer \emph{where was the object observed?}; navigation requires \emph{where is it likely to be now?} More importantly, a predictive answer should not be consumed as a static waypoint. Travelling takes time, candidate locations differ in inspection cost and visibility, and a failed inspection is new evidence. A useful memory must therefore support a closed loop: retrieve relevant history, propagate it to the query time, select a cost-effective location, observe, update the belief, and replan.

Several recent lines of work approach parts of this problem. Dynamic scene memories update object states after new observations \citep{liu2024dynamem,kurenkov2023dynamic}; predictive scene graphs estimate the persistence or re-emergence of time-varying relations \citep{saavedraruiz2026predictivegraphs}; and models of household routines predict multimodal future object placements for downstream navigation \citep{patel2023homer,argenziano2026flowmaps}. Portable Object Navigation goes further by allowing a target to move while the agent travels and learning transit-aware policies \citep{dorbala2026portable}. These methods establish that dynamics matter, but leave a complementary question open: how should an agent turn its own long-term, partially observed memory into a belief that directly controls physical search and is corrected by the evidence collected along the route?

We formulate this setting as \textbf{persistent dynamic object navigation}. Before receiving a query, the agent has accumulated a timestamped 4D memory through intermittent visits. The target may undergo one or more hidden transitions after its last positive observation. At query time, the agent receives neither the current state nor the activity that caused the transition. It must find the object under a navigation budget using only causal memory, a map of candidate receptacles and viewpoints, and online RGB-D observations. Our main protocol freezes the target once the query begins to isolate the value of predictive memory; an extension allows an additional hidden transition during navigation.

We introduce \pfdnav{}, a compact predictive-memory agent for this problem. Its 4D memory maintains persistent entity identities, time-valid state versions, and observation provenance in a common 3D frame. A target-conditioned retriever exposes only the history available before the query. A lightweight continuous-time Transformer represents irregular temporal gaps and calendar periodicity. Instead of directly classifying a location, a structured head factorizes the belief into (i) the probability that the last-seen state persists and (ii) a pointer distribution over relocation candidates. This inductive bias matches the physical alternatives faced by a navigator and supports candidate sets that vary across scenes.

Prediction is coupled to action through a visibility-aware evidence update and a belief-guided controller. Concentrated beliefs support a direct first attempt; diffuse beliefs induce a probability--cost search order; and sufficiently covered negative observations suppress inspected hypotheses before replanning. The controller is identical across stable, routine-mobile, personal-habit, and irregular objects and receives no mobility label at test time. Thus, static recall, predictive navigation, and uncertainty-driven search emerge from one belief rather than from an oracle router.

To evaluate the resulting navigation problem, we introduce \pfdbench{}. Its dynamic HSSD episodes are compiled through a canonical routine representation that combines complementary signals from longitudinal activity records, real human--object interactions, and household object-state scripts. More importantly, the benchmark contains paired static, routine, and random worlds that share the scene, target, last observation, query time, candidate set, and agent start. Only the hidden transition process differs. This counterfactual construction tests whether an agent exploits temporal structure instead of object--location frequency or route-length shortcuts.

\paragraph{Contributions.}
Our contributions are threefold:
\begin{itemize}[leftmargin=5mm,itemsep=0.6mm,topsep=1mm]
    \item We formulate persistent dynamic object navigation, connecting long-term, partially observed memory to closed-loop physical search under hidden object transitions.
    \item We propose \pfdnav{}, which combines an evidence-grounded 4D memory, continuous-time history encoding, a persistence--relocation belief, and visibility-aware Bayesian correction to support direct navigation, belief-guided search, and recovery without oracle dynamics labels.
    \item We introduce \pfdbench{}, a human-trace-grounded HSSD benchmark with paired dynamic controls, instance-level mobility regimes, online replanning tasks, and cross-scene and cross-household evaluation.
\end{itemize}
